\chapter{The Infinite Galois Correspondence}
\label{ch:galois-correspondence}

Throughout Sections~\ref{sec:fixed-fields}--\ref{sec:normal-quotients},
let \(L/K\) be Galois, let \(G=\operatorname{Gal}(L/K)\) carry the Krull
topology, and let \(G_M=\operatorname{Gal}(L/M)\) for an intermediate
field \(M\). For a subgroup \(H\leq G\) we write
\[
  L^H=\{x\in L : h(x)=x\text{ for every }h\in H\}.
\]
This is a field between \(K\) and \(L\), since every \(h\) preserves
addition, multiplication, and inverses of nonzero elements. We keep
the notation \(Q_E=\operatorname{Gal}(E/K)\) and
\(r_E:G\to Q_E\) for a finite Galois intermediate field \(E\).
The finite fixed-field theorem applied in \(E\) will control both
the restriction of \(L^H\) to \(E\) and the closure of \(H\) in \(G\).

\section{Fixed fields and closure}
\label{sec:fixed-fields}

\begin{proposition}[The fixed field of the full group]
\label{res:C01}
For a Galois extension \(L/K\), the full group
\(G=\operatorname{Gal}(L/K)\) has fixed field \(K\). More precisely,
if \(x\in L\setminus K\), then there is a finite Galois intermediate
field \(E\) containing \(x\) and an element of \(Q_E\) moving \(x\)
that is the restriction of an element of \(G\).
\end{proposition}

\begin{proof}
Choose \(E\) by \ref{res:F03}. The finite fixed-field theorem
of \ref{res:F01} gives \(E^{Q_E}=K\), so some \(\tau\in Q_E\)
satisfies \(\tau(x)\ne x\). Surjectivity of \(r_E\) in \ref{res:F07}
provides \(g\in G\) with \(g|_E=\tau\). Therefore \(g(x)\ne x\).
All elements of \(K\) are fixed by the definition of \(G\).
\cite[Proposition 7.9]{milne}
\end{proof}

\begin{proposition}[Recovery of an intermediate field]
\label{res:C02}
For every intermediate field \(M\), the subgroup \(G_M\) is closed
and \(L^{G_M}=M\).
\end{proposition}

\begin{proof}
For \(x\in M\), the stabilizer of \(x\) in \(G\) is the inverse
image of \(\{x\}\) under evaluation at \(x\). Evaluation is continuous
to the discrete space \(L\) by \ref{res:T03}, so the stabilizer is
closed. Their intersection over \(x\in M\) is \(G_M\). By \ref{res:F02}
the extension \(L/M\) is Galois, and \ref{res:C01} with base field \(M\)
gives \(L^{G_M}=M\).
\cite[Proposition 7.12(a)]{milne}
\end{proof}

\begin{lemma}[Fixed fields on a finite layer]
\label{res:C03}
For any subgroup \(H\leq G\) and any finite Galois intermediate field
\(E\), one has
\[
  L^H\cap E=E^{r_E(H)}
  \qquad\text{and}\qquad
  r_E(H)=\operatorname{Gal}\bigl(E/(L^H\cap E)\bigr).
\]
\end{lemma}

\begin{proof}
Each element \(h\in H\) preserves \(E\) because \(E/K\) is normal.
For \(x\in E\) the condition \(h(x)=x\) for every \(h\in H\)
is therefore the condition that each restricted automorphism in
\(r_E(H)\) fix \(x\). This proves the first equality. The image
\(r_E(H)\) is a subgroup of the finite Galois group \(Q_E\).
Apply the finite correspondence of \ref{res:F01} to that subgroup
and the intermediate field \(E^{r_E(H)}\) to obtain the second equality.
\cite[Tag 0BML]{stacks}
\end{proof}

\begin{theorem}[Closure and the infinite Galois correspondence]
\label{res:C04}
For any subgroup \(H\leq G\),
\[
    G_{L^H}=\overline H,
    \qquad L^{\overline H}=L^H .
\]
Consequently the assignments \(M\mapsto G_M\) and \(H\mapsto L^H\)
give inverse inclusion-reversing bijections between the intermediate
fields of \(L/K\) and the closed subgroups of \(G\).
\end{theorem}

\begin{proof}
The group \(G_{L^H}\) is closed by \ref{res:C02} and contains \(H\),
so it contains \(\overline H\). If \(g\in G_{L^H}\), then \(r_E(g)\)
fixes every element of \(L^H\cap E\). By \ref{res:C03} this says that
\(r_E(g)\in r_E(H)\). This is true for every finite Galois intermediate
field \(E\), so the finite-projection closure criterion \ref{res:T07}
gives \(g\in\overline H\). We have proved the first equality. In
particular \(\overline H\) is a subgroup, and \ref{res:C02} applied
to \(M=L^H\) gives the second equality.

For a closed subgroup \(H\), the first equality now reads
\(G_{L^H}=H\). This and \ref{res:C02} prove that the assignments
are inverse. Enlarging an intermediate field imposes more pointwise
fixing conditions on its subgroup. Enlarging a subgroup imposes
more conditions on its fixed field. Thus the two inverses reverse
inclusion.
\cite[Proposition 7.12 and Theorem 7.13(a)]{milne}
\cite[Theorem 2.11.3]{rz}
\cite[Tag 0BML]{stacks}
\end{proof}

\begin{corollary}[Lattice operations and conjugation]
\label{res:C05}
For a family of intermediate fields \((M_i)_{i\in I}\), let
\(\bigvee_i M_i\) be their compositum in \(L\), and put \(H_i=G_{M_i}\).
Then
\[
  G_{\bigvee_i M_i}=\bigcap_i H_i,
  \qquad
  G_{\bigcap_i M_i}
     =\overline{\left\langle\bigcup_i H_i\right\rangle}.
\]
An empty compositum is taken to be \(K\), and an empty intersection
of intermediate fields is \(L\). For \(g\in G\) one also has
\[
  G_{g(M)}=gG_Mg^{-1},
  \qquad
  L^{gHg^{-1}}=g(L^H)\quad\text{for every }H\leq G.
\]
\end{corollary}

\begin{proof}
Fixing all \(M_i\) pointwise is equivalent to fixing the field they
generate pointwise, which proves the first lattice equality. Write
\(J=\langle\bigcup_i H_i\rangle\). An element of \(L\) is fixed by \(J\)
exactly when it is fixed by all the \(H_i\), because a product of
automorphisms fixing an element and their inverses still fixes the
element. Thus \(L^J=\bigcap_i L^{H_i}=\bigcap_i M_i\) by
\ref{res:C02}. Apply \ref{res:C04} to \(J\).

An automorphism \(a\in G\) fixes all elements \(g(m)\), \(m\in M\),
exactly when \(g^{-1}ag\) fixes every \(m\in M\). This proves the
first conjugation formula. For the second, write a candidate element
as \(g(x)\). Since \(ghg^{-1}(g(x))=g(h(x))\), it is fixed by all
the conjugates of \(H\) exactly when \(x\) is fixed by \(H\).
\cite[Theorem 7.13(c) and Remark 7.14]{milne}
\end{proof}

\section{Finite intermediate fields and cosets}

\begin{theorem}[Open subgroups and finite extensions]
\label{res:C06}
For an intermediate field \(M\), the subgroup \(G_M\) is open if
and only if \(M/K\) is finite. More generally, a closed subgroup
of \(G\) is open if and only if it has finite index.
\end{theorem}

\begin{proof}
If \(M/K\) is finite, a finite \(K\)-basis \(S\) for \(M\) generates
\(M\) as a field over \(K\). The subgroup \(G_M\) is then the
pointwise stabilizer of \(S\), which is open by \ref{res:T01}.
Conversely, if \(G_M\) is open, it contains the subgroup \(N_E=G_E\)
for some finite Galois intermediate field \(E\), by the neighborhood
basis of \ref{res:T02}. Taking fixed fields and applying \ref{res:C02}
gives \(M=L^{G_M}\subseteq L^{G_E}=E\). The inclusion makes \(M\)
a \(K\)-subspace of the finite-dimensional space \(E\), so \(M/K\)
is finite.

If \(H\leq G\) is open, its left cosets are open and form a disjoint
cover of the compact space \(G\) by \ref{res:T06}. A finite subcover
must contain all cosets, so the index is finite. If \(H\) is closed
of finite index, its complement is the finite union of its other
left cosets. Each is closed because left translation is a
homeomorphism. Thus \(H\) is open.
\cite[Theorem 7.13(b)]{milne}
\cite[Sections 2.1 and 2.11]{rz}
\end{proof}

\begin{proposition}[Cosets as embeddings]
\label{res:C07}
The restriction of automorphisms induces a bijection of sets
\[
  G/G_M \longrightarrow \operatorname{Emb}_K(M,L),
  \qquad gG_M\longmapsto g|_M ,
\]
where \(G/G_M\) denotes the set of left cosets. In particular, if
\(M/K\) is finite, then
\[
    [G:G_M]=[M:K].
\]
\end{proposition}

\begin{proof}
The equality \(g|_M=g'|_M\) holds if and only if
\(g^{-1}g'\in G_M\), so the map is well-defined and injective.
Every \(K\)-embedding \(M\to L\) extends to an element of \(G\)
by \ref{res:F05}, giving surjectivity. For finite \(M/K\), it is
separable by \ref{res:F02}. Choose a primitive element \(x\) by
the finite preliminaries \ref{res:F01}. Its minimal polynomial
over \(K\) has degree \([M:K]\), has distinct roots, and splits
in \(L\) by normality. Each root gives exactly one \(K\)-embedding
of \(K(x)=M\) into \(L\), proving the count. This identifies the
degree with the coset index for every finite intermediate field,
including those for which \(G_M\) is not normal.
\cite[Proposition 7.4 and Theorem 7.13(b)]{milne}
\cite[Chapters 12 and 14]{d07}
\end{proof}

\section{Normal subgroups and topological quotients}
\label{sec:normal-quotients}

\begin{corollary}[Normality in the correspondence]
\label{res:C08}
Let \(H\) be a closed subgroup of \(G\) with fixed field \(M\).
Then \(H\) is normal in \(G\) if and only if \(M/K\) is Galois.
The open normal subgroups correspond precisely to the finite
Galois intermediate extensions of \(K\).
\end{corollary}

\begin{proof}
By \ref{res:C04}, \(H=G_M\). For every \(g\in G\), the conjugate
\(gG_Mg^{-1}\) equals \(G_{g(M)}\) by \ref{res:C05}.
By injectivity of the field-to-subgroup correspondence,
this equals \(G_M\) exactly when \(g(M)=M\). By \ref{res:F06},
stability under all \(g\) is equivalent to the Galois condition
on \(M/K\). An open subgroup is closed, since its complement is a
union of open cosets, so the openness condition can be added
using \ref{res:C06}.
\cite[Theorem 7.13(c)--(d)]{milne}
\cite[Tag 0BML]{stacks}
\end{proof}

\begin{theorem}[The profinite exact sequence]
\label{res:C09}
Suppose \(M/K\) is Galois. The inclusion and restriction maps form
a short exact sequence of profinite topological groups
\[
  1\longrightarrow G_M\longrightarrow G
   \xrightarrow{\ r_M\ }\operatorname{Gal}(M/K)
   \longrightarrow 1 .
\]
In particular the induced map
\[
    G/G_M \longrightarrow \operatorname{Gal}(M/K),
    \qquad gG_M\longmapsto g|_M ,
\]
is a topological group isomorphism for the quotient topology on
the left and the Krull topology on the right.
\end{theorem}

\begin{proof}
The subgroup \(G_M\) is closed and normal by \ref{res:C02}
and \ref{res:C08}. Its own Krull topology as \(\operatorname{Gal}(L/M)\)
is exactly the topology induced by its inclusion into \(G\), by
\ref{res:T08}. The restriction \(r_M\) is continuous by \ref{res:T03},
and \ref{res:F07} identifies its kernel as \(G_M\) and proves
surjectivity. Hence it factors as a bijective homomorphism through
the group quotient. The factor is continuous by the definition
of the quotient topology. Its domain is compact as a continuous
image of \(G\), and \(\operatorname{Gal}(M/K)\) is Hausdorff by
\ref{res:T06}. A continuous bijection from a compact space to
a Hausdorff space is a homeomorphism, since it sends closed sets
to compact and hence closed sets. The inclusion is a topological
embedding, and the restriction is a quotient map by this homeomorphism.
All three groups are profinite by \ref{res:T06}, applied to
the Galois extensions \(L/M\), \(L/K\), and \(M/K\).
\cite[Theorem 7.13(d)]{milne}
\cite[Theorem 2.11.3]{rz}
\cite[Tag 0BMM]{stacks}
\end{proof}

\section{Restriction after a base change}

\begin{theorem}[Galois base change in a common overfield]
\label{res:C10}
Let \(E\) and \(M\) be extensions of \(K\) in a common overfield,
and suppose \(E/K\) is Galois. Then \(EM/M\) and \(E/(E\cap M)\)
are Galois, and restriction is an isomorphism of topological groups
\[
  \operatorname{Gal}(EM/M)
     \longrightarrow \operatorname{Gal}\bigl(E/(E\cap M)\bigr),
     \qquad g\longmapsto g|_E .
\]
\end{theorem}

\begin{proof}
The extension \(E/(E\cap M)\) is Galois by \ref{res:F02}.
By \ref{res:F04}, \(E\) is the directed union of its finite
Galois intermediate fields \(E_i\) over \(K\). Each \(E_i\)
is the splitting field of a finite collection of separable
polynomials over \(K\): take minimal polynomials of a finite
set of field generators and use normality to put all their
roots in \(E_i\). Over \(M\) the same polynomials remain separable,
and their roots generate \(E_iM\) over \(M\). Hence \(E_iM/M\)
is finite Galois. The union of the \(E_iM\) is a field containing
both \(E\) and \(M\), and it is contained in \(EM\), so it equals
\(EM\). The converse direction of \ref{res:F04} now proves
that \(EM/M\) is Galois.

Every \(g\in\operatorname{Gal}(EM/M)\) is a \(K\)-automorphism and
sends \(E\) into \(E\) by normality of \(E/K\): for \(x\in E\),
the element \(g(x)\) is a root of the \(K\)-minimal polynomial of
\(x\), and all those roots already belong to \(E\). The same applies
to \(g^{-1}\), so \(g|_E\) is an automorphism fixing \(E\cap M\).
The restriction homomorphism is injective because \(E\) and \(M\)
generate \(EM\). For a finite subset \(S\subseteq E\), the inverse
image of its pointwise stabilizer in the target is its pointwise
stabilizer in the source, which is open by \ref{res:T01}.
Thus the homomorphism is continuous at the identity and hence
everywhere.

Let \(H\) denote its image. The source is compact and the target
is Hausdorff by \ref{res:T06}, so \(H\) is closed. An element of
\(E\) is fixed by \(H\) exactly when it is fixed by every
element of \(\operatorname{Gal}(EM/M)\). By \ref{res:C01},
the latter group fixes precisely \(M\) in \(EM\). It follows that
\(E^H=E\cap M\). Apply \ref{res:C04} to the Galois extension
\(E/(E\cap M)\) and its closed subgroup \(H\) to obtain
\(H=\operatorname{Gal}(E/(E\cap M))\). Restriction is now a
continuous bijection from compact to Hausdorff as in \ref{res:C09},
and therefore is a homeomorphism.
\cite[Proposition 7.15]{milne}
\cite[Chapter 14, finite base change]{d07}
\end{proof}

\section{Fixed fields of compact groups and realization}

\begin{theorem}[Fixed fields of faithful compact actions]
\label{res:C11}
Let \(G\) be a compact topological group acting faithfully by field
automorphisms on a field \(E\). Suppose the stabilizer of every
element of \(E\) is open in \(G\). Put \(K=E^G\). Then \(E/K\)
is Galois, and the homomorphism induced by the action
\[
 G \longrightarrow \operatorname{Gal}(E/K)
\]
is an isomorphism of topological groups, with the Krull topology
on the right.
\end{theorem}

\begin{proof}
The orbit of \(a\in E\) is in bijection with the left cosets of its
stabilizer. These cosets are open and cover the compact space \(G\).
A finite subcover therefore accounts for every coset, so every
orbit is finite.

For a finite set \(S\subseteq E\), let \(E_S\) be the field generated
over \(K\) by all the \(G\)-orbits of elements of \(S\). The
generator set is finite and \(G\)-stable. Hence \(G\) acts on
\(E_S\) through a finite group \(H_S\) of field automorphisms:
the action on the generator set determines the restriction to
the whole field. The fixed field of \(H_S\) in \(E_S\) is exactly
\(K\), since an element of \(E_S\) fixed by \(H_S\) is fixed by
all of \(G\). Artin's finite fixed-field theorem, in
\cite[Chapter 14, Artin's fixed-field theorem]{d07} and
\cite[Corollary 3.11]{milne}, now gives that \(E_S/K\) is
finite Galois and that \(H_S=\operatorname{Gal}(E_S/K)\).
In particular the restricted action of \(G\) fills this entire
finite Galois group.

The fields \(E_S\) are directed since \(E_{S\cup T}\) contains
both \(E_S\) and \(E_T\), and their union is \(E\). The converse
in \ref{res:F04} makes \(E/K\) Galois. Any finite Galois
intermediate field is generated by finitely many elements over
\(K\), so it is contained in some \(E_S\).

The homomorphism \(G\to\operatorname{Gal}(E/K)\) is injective
by faithfulness. It is continuous because the inverse image of
a basic pointwise stabilizer on a finite set is an intersection
of finitely many open stabilizers in \(G\), using
\ref{res:T01}. Its image is dense: given an automorphism of
\(E/K\) and any finite set \(S\) on which to match it, take its
restriction to \(E_S\). The equality
\(H_S=\operatorname{Gal}(E_S/K)\) provides an element of \(G\)
with that restriction. This meets each Krull neighborhood of the
given automorphism. The image is also compact, hence closed in
the Hausdorff group of \ref{res:T06}. It is therefore the whole
group. A continuous bijection from compact to Hausdorff is a
homeomorphism, as in \ref{res:C09}.
\cite[Proposition 7.10]{milne}
\end{proof}

\Needspace{16\baselineskip}
\begin{corollary}[Tate's realization theorem]
\label{res:C12}
Every profinite group \(G\) is isomorphic, as a topological
group, to \(\operatorname{Gal}(E/K)\) for some Galois extension
\(E/K\). The fields can be constructed in any specified
characteristic.
\end{corollary}

\begin{proof}
Let \(X\) be the disjoint union of the sets \(G/U\) over the open
normal subgroups \(U\) of \(G\). Such subgroups have finite
index by compactness. The left translation action on \(X\) is
faithful: the kernel on \(G/U\) is \(U\), and the intersection
of the open normal subgroups is trivial. Indeed, in an inverse
limit presentation of \(G\), the finite-coordinate kernels
are open and normal and separate points. Every point of \(X\)
has open stabilizer, because \(U\) fixes each point of the
corresponding component.

Choose a field \(k\) of the desired characteristic and put
\(E=k(t_x:x\in X)\), the rational function field in independent
indeterminates indexed by \(X\). The rule
\(t_x\mapsto t_{g x}\) permutes the variables and extends
to a \(k\)-automorphism for each \(g\in G\). These
automorphisms form a faithful action. A rational function
involves only finitely many variables. Its stabilizer contains
the intersection of their point stabilizers, an open subgroup,
and is therefore open as a union of cosets of that subgroup.
Apply \ref{res:C11} to this action, with \(K=E^G\).
The construction specifies \(K\) as a fixed field of \(E\).
\cite[Proposition 7.27]{milne}
\end{proof}
